\begin{longtable}{|p{3.5cm}|p{5.5cm}|p{6cm}|}

\hline

Categoria & Métricas identificadas na literatura & Exemplos de estudos \\ \hline

\endfirsthead

\hline

Categoria & Métricas identificadas na literatura & Exemplos de estudos \\ \hline

\endhead

\hline

\endfoot

\endlastfoot

Desempenho na realização de tarefas &

Tempo de tarefa, tempo de conclusão, tempo de busca, taxa de conclusão de tarefas, taxa de sucesso, número de cliques, facilidade de encontrar informações, qualidade da informação, precisão, acurácia, recomendação e desempenho durante a realização das tarefas. &

Development of a Weighted Heuristic for Website Evaluation for Older Adults; Development of a web-based insulin decision aid for the elderly; Usability Testing by Older Americans of a Prototype Google Map Web Site To Select Nursing Homes. \\ \hline

Usabilidade percebida e satisfação &

System Usability Scale (SUS), PSSUQ, e-Qual, escalas Likert de satisfação, facilidade de uso, percepção de usabilidade, qualidade percebida do design, satisfação do usuário e experiência percebida. &

Development of a Weighted Heuristic for Website Evaluation for Older Adults; Usability Assessment of Secure Messaging for Clinical Document Sharing between Health Care Providers and Patients; Novice user perception of e-services: A study in the Egyptian public sector. \\ \hline

Eficiência, eficácia e navegação &

Eficácia, eficiência, taxa de sucesso, tempo de conclusão, tempo de execução, path length, category node hits, frequência de uso declarada e desempenho na recuperação de informações. &

Why structure and genre matter for users of digital information: A longitudinal experiment with readers of a web-based newspaper; Age related differences and the depth vs. breadth tradeoff in hierarchical online information systems; Usability evaluation of a social networking site prototype for the elderly. \\ \hline

Erros e problemas de interação &

Quantidade de erros, tipos de erros, erros de uso, número de problemas identificados, severidade dos problemas, complexidade, problemas críticos, sérios ou cosméticos, capacidade de recuperação de erros e não-problemas identificados durante a avaliação. &

Usability Assessment of Secure Messaging for Clinical Document Sharing between Health Care Providers and Patients; A conceptual tool for usability problem identification in website development; Using an Error Detection Strategy for Improving Web Accessibility for Older Adults. \\ \hline

Acessibilidade e conformidade &

Conformidade com WCAG 2.0, conformidade com WCAG 2.0 AA, violações de acessibilidade, resultados de ferramentas de avaliação automática, resultados do WAVE, resultados do A-Checker, validação de marcação e validação CSS. &

A methodological approach for designing and developing web-based inventories of mobile Assistive Technology applications; Family-centered design: Interactive performance testing and user interface evaluation of the Slovenian eDavki public tax portal; Guideline for Web pages design inclusive for senior users. \\ \hline

Experiência, percepção e fatores subjetivos &

Ansiedade, satisfação, facilidade de uso percebida, desorientação percebida, confiança, autoeficácia, literacia em IA, qualidade percebida do design e percepção da experiência de uso. &

Exploring the wiki user experience: The effects of training spaces on novice user usability and anxiety towards wiki editing; Age, technology usage, and cognitive characteristics in relation to perceived disorientation and reported website ease of use; U-Cker: Initial Development of an Interactive Video Retrieval System for Novice Users. \\ \hline

Aspectos visuais, cognitivos e de interação &

Contagem de fixações, contagem de sacadas, tempo até a primeira fixação, diâmetro pupilar, quadtree leaves, clutter, balance, equilibrium, número de imagens, número de fontes, compreensão, recordação, reconhecimento, complexidade de leitura, ARI e Kincaid. &

Age-Related Differences in Eye Tracking and Usability Performance: Website Usability for Older Adults; HCI for Elderly Measuring Visual Complexity of Webpages Based on Machine Learning; Why structure and genre matter for users of digital information: A longitudinal experiment with readers of a web-based newspaper. \\ \hline

Métricas específicas da tarefa ou do contexto &

Insights por minuto, tipos de insight, correção dos insights, path length, category node hits, tempo de download, frequência de uso, aceitação tecnológica e medidas relacionadas à realização de tarefas específicas. &

Keshif: Rapid and Expressive Tabular Data Exploration for Novices; Why structure and genre matter for users of digital information: A longitudinal experiment with readers of a web-based newspaper; A study of web usability for older adults seeking online health resources. \\ \hline

\caption{Métricas de usabilidade e acessibilidade identificadas nos estudos que respondem diretamente à RQ1.}

\label{tab:metricasRQ1}

\end{longtable}